\documentclass[reqno]{amsart}
\usepackage[T1]{fontenc}
\usepackage{amsmath,amssymb,amsthm}
\usepackage{xr-hyper}
\usepackage[hidelinks]{hyperref}
\externaldocument[da-]{20261126T030300--the-data__workbook}
\externaldocument[ob-]{20261126T040100--spaces-forms-and-twists__workbook}
\externaldocument[ho-]{20261125T016500--hecke-operators-and-theta-lifts__workbook}
\externaldocument[dch-]{20261126T030700--the-determinant-of-a-chain__workbook}
\externaldocument[c1-]{20261126T050400--the-norm-of-one-lift__workbook}

% Chapter number; \ref{chap} gives it.
\newcounter{chap}
% Chapter numbers, generated from 20261125T012000--grh-workbook__guide.org
% by scripts/make-chapters.py.  Do not edit: change the order of the
% index and rerun the script (./preview-site.sh and CI do this).
%
% Usage in a sheet, after \newcounter{chap}:
%   % Chapter numbers, generated from 20261125T012000--grh-workbook__guide.org
% by scripts/make-chapters.py.  Do not edit: change the order of the
% index and rerun the script (./preview-site.sh and CI do this).
%
% Usage in a sheet, after \newcounter{chap}:
%   % Chapter numbers, generated from 20261125T012000--grh-workbook__guide.org
% by scripts/make-chapters.py.  Do not edit: change the order of the
% index and rerun the script (./preview-site.sh and CI do this).
%
% Usage in a sheet, after \newcounter{chap}:
%   \input{chapters}\chapnumber{<sheet ID>}
\makeatletter
\providecommand\chapters@set[2]{\@namedef{chapters@#1}{#2}}
\chapters@set{20261125T001000}{1}% zeros-to-ramp-sums
\chapters@set{20261125T001100}{2}% inert-primes-and-estimate-b
\chapters@set{20261125T005000}{3}% galois-fourier-and-averaging
\chapters@set{20261125T005500}{4}% columns-from-a-sphere
\chapters@set{20261125T000000}{5}% slope-method
\chapters@set{20261126T030100}{6}% a-twist-for-each-prime
\chapters@set{20261125T003000}{7}% the-chain-and-the-square-root
\chapters@set{20261125T010000}{8}% rank-lemma-and-zero-estimate
\chapters@set{20261125T010500}{9}% poincare-series-and-near-orthogonality
\chapters@set{20261124T223000}{10}% elliptic-curves-for-the-proof
\chapters@set{20261125T009500}{11}% fock-model-and-theta-intertwiners
\chapters@set{20261124T223100}{12}% complex-multiplication-and-inert-reduction
\chapters@set{20261125T008000}{13}% one-cm-determinant
\chapters@set{20261124T214500}{14}% supersingular-theta-divisibility
\chapters@set{20261124T214600}{15}% the-isotropic-plane
\chapters@set{20261125T013700}{16}% pairings-on-a-formal-group-of-height-two
\chapters@set{20261125T013500}{17}% hecke-nondegeneracy-on-the-tree
\chapters@set{20261126T030200}{18}% the-bound-is-attained
\chapters@set{20261125T009000}{19}% kronecker-limit-and-chowla-selberg
\chapters@set{20261125T002000}{20}% definite-theta-lifts
\chapters@set{20261125T015500}{21}% siegel-weil-in-the-smallest-case
\chapters@set{20261125T002100}{22}% definite-theta-lifts-and-norms
\chapters@set{20261125T016500}{23}% hecke-operators-and-theta-lifts
\chapters@set{20261125T016700}{24}% divisibility-along-the-chain
\chapters@set{20261125T004000}{25}% denominators-and-sha
\chapters@set{20261125T004100}{26}% tate-twists-and-sha
\chapters@set{20261126T020430}{27}% the-minus-part-by-gauss-sums
\chapters@set{20261125T014100}{28}% congruence-modules
\chapters@set{20261125T014500}{29}% eisenstein-congruences-and-selmer-classes
\chapters@set{20261125T014700}{30}% pseudocharacters-and-the-length-bound
\chapters@set{20261125T011000}{31}% galois-representations-black-box
\chapters@set{20261126T030300}{32}% the-data
\chapters@set{20261126T020210}{33}% bessel-poincare-series-and-their-gram-operators
\chapters@set{20261126T030400}{34}% the-seesaw-identity
\chapters@set{20261126T030500}{35}% divisibility-of-the-composite-lift
\chapters@set{20261126T030600}{36}% theorem-a-short-forms
\chapters@set{20261126T030700}{37}% the-determinant-of-a-chain
\chapters@set{20261126T030800}{38}% the-cost-of-one-lift
\chapters@set{20261126T030900}{39}% denominators-of-projections-to-lifts
\chapters@set{20261126T031000}{40}% twists-with-few-congruences
\chapters@set{20261126T031100}{41}% theorem-c-and-the-main-inequality
\providecommand\chapnumber[1]{%
  \@ifundefined{chapters@#1}%
    {\PackageError{chapters}{Sheet #1 is not listed in the index}%
      {List it in 20261125T012000--grh-workbook__guide.org\MessageBreak
       and run scripts/make-chapters.py.}}%
    {\edef\thechap{\@nameuse{chapters@#1}}}}
\makeatother
\chapnumber{<sheet ID>}
\makeatletter
\providecommand\chapters@set[2]{\@namedef{chapters@#1}{#2}}
\chapters@set{20261125T001000}{1}% zeros-to-ramp-sums
\chapters@set{20261125T001100}{2}% inert-primes-and-estimate-b
\chapters@set{20261125T005000}{3}% galois-fourier-and-averaging
\chapters@set{20261125T005500}{4}% columns-from-a-sphere
\chapters@set{20261125T000000}{5}% slope-method
\chapters@set{20261126T030100}{6}% a-twist-for-each-prime
\chapters@set{20261125T003000}{7}% the-chain-and-the-square-root
\chapters@set{20261125T010000}{8}% rank-lemma-and-zero-estimate
\chapters@set{20261125T010500}{9}% poincare-series-and-near-orthogonality
\chapters@set{20261124T223000}{10}% elliptic-curves-for-the-proof
\chapters@set{20261125T009500}{11}% fock-model-and-theta-intertwiners
\chapters@set{20261124T223100}{12}% complex-multiplication-and-inert-reduction
\chapters@set{20261125T008000}{13}% one-cm-determinant
\chapters@set{20261124T214500}{14}% supersingular-theta-divisibility
\chapters@set{20261124T214600}{15}% the-isotropic-plane
\chapters@set{20261125T013700}{16}% pairings-on-a-formal-group-of-height-two
\chapters@set{20261125T013500}{17}% hecke-nondegeneracy-on-the-tree
\chapters@set{20261126T030200}{18}% the-bound-is-attained
\chapters@set{20261125T009000}{19}% kronecker-limit-and-chowla-selberg
\chapters@set{20261125T002000}{20}% definite-theta-lifts
\chapters@set{20261125T015500}{21}% siegel-weil-in-the-smallest-case
\chapters@set{20261125T002100}{22}% definite-theta-lifts-and-norms
\chapters@set{20261125T016500}{23}% hecke-operators-and-theta-lifts
\chapters@set{20261125T016700}{24}% divisibility-along-the-chain
\chapters@set{20261125T004000}{25}% denominators-and-sha
\chapters@set{20261125T004100}{26}% tate-twists-and-sha
\chapters@set{20261126T020430}{27}% the-minus-part-by-gauss-sums
\chapters@set{20261125T014100}{28}% congruence-modules
\chapters@set{20261125T014500}{29}% eisenstein-congruences-and-selmer-classes
\chapters@set{20261125T014700}{30}% pseudocharacters-and-the-length-bound
\chapters@set{20261125T011000}{31}% galois-representations-black-box
\chapters@set{20261126T030300}{32}% the-data
\chapters@set{20261126T020210}{33}% bessel-poincare-series-and-their-gram-operators
\chapters@set{20261126T030400}{34}% the-seesaw-identity
\chapters@set{20261126T030500}{35}% divisibility-of-the-composite-lift
\chapters@set{20261126T030600}{36}% theorem-a-short-forms
\chapters@set{20261126T030700}{37}% the-determinant-of-a-chain
\chapters@set{20261126T030800}{38}% the-cost-of-one-lift
\chapters@set{20261126T030900}{39}% denominators-of-projections-to-lifts
\chapters@set{20261126T031000}{40}% twists-with-few-congruences
\chapters@set{20261126T031100}{41}% theorem-c-and-the-main-inequality
\providecommand\chapnumber[1]{%
  \@ifundefined{chapters@#1}%
    {\PackageError{chapters}{Sheet #1 is not listed in the index}%
      {List it in 20261125T012000--grh-workbook__guide.org\MessageBreak
       and run scripts/make-chapters.py.}}%
    {\edef\thechap{\@nameuse{chapters@#1}}}}
\makeatother
\chapnumber{<sheet ID>}
\makeatletter
\providecommand\chapters@set[2]{\@namedef{chapters@#1}{#2}}
\chapters@set{20261125T001000}{1}% zeros-to-ramp-sums
\chapters@set{20261125T001100}{2}% inert-primes-and-estimate-b
\chapters@set{20261125T005000}{3}% galois-fourier-and-averaging
\chapters@set{20261125T005500}{4}% columns-from-a-sphere
\chapters@set{20261125T000000}{5}% slope-method
\chapters@set{20261126T030100}{6}% a-twist-for-each-prime
\chapters@set{20261125T003000}{7}% the-chain-and-the-square-root
\chapters@set{20261125T010000}{8}% rank-lemma-and-zero-estimate
\chapters@set{20261125T010500}{9}% poincare-series-and-near-orthogonality
\chapters@set{20261124T223000}{10}% elliptic-curves-for-the-proof
\chapters@set{20261125T009500}{11}% fock-model-and-theta-intertwiners
\chapters@set{20261124T223100}{12}% complex-multiplication-and-inert-reduction
\chapters@set{20261125T008000}{13}% one-cm-determinant
\chapters@set{20261124T214500}{14}% supersingular-theta-divisibility
\chapters@set{20261124T214600}{15}% the-isotropic-plane
\chapters@set{20261125T013700}{16}% pairings-on-a-formal-group-of-height-two
\chapters@set{20261125T013500}{17}% hecke-nondegeneracy-on-the-tree
\chapters@set{20261126T030200}{18}% the-bound-is-attained
\chapters@set{20261125T009000}{19}% kronecker-limit-and-chowla-selberg
\chapters@set{20261125T002000}{20}% definite-theta-lifts
\chapters@set{20261125T015500}{21}% siegel-weil-in-the-smallest-case
\chapters@set{20261125T002100}{22}% definite-theta-lifts-and-norms
\chapters@set{20261125T016500}{23}% hecke-operators-and-theta-lifts
\chapters@set{20261125T016700}{24}% divisibility-along-the-chain
\chapters@set{20261125T004000}{25}% denominators-and-sha
\chapters@set{20261125T004100}{26}% tate-twists-and-sha
\chapters@set{20261126T020430}{27}% the-minus-part-by-gauss-sums
\chapters@set{20261125T014100}{28}% congruence-modules
\chapters@set{20261125T014500}{29}% eisenstein-congruences-and-selmer-classes
\chapters@set{20261125T014700}{30}% pseudocharacters-and-the-length-bound
\chapters@set{20261125T011000}{31}% galois-representations-black-box
\chapters@set{20261126T030300}{32}% the-data
\chapters@set{20261126T020210}{33}% bessel-poincare-series-and-their-gram-operators
\chapters@set{20261126T030400}{34}% the-seesaw-identity
\chapters@set{20261126T030500}{35}% divisibility-of-the-composite-lift
\chapters@set{20261126T030600}{36}% theorem-a-short-forms
\chapters@set{20261126T030700}{37}% the-determinant-of-a-chain
\chapters@set{20261126T030800}{38}% the-cost-of-one-lift
\chapters@set{20261126T030900}{39}% denominators-of-projections-to-lifts
\chapters@set{20261126T031000}{40}% twists-with-few-congruences
\chapters@set{20261126T031100}{41}% theorem-c-and-the-main-inequality
\providecommand\chapnumber[1]{%
  \@ifundefined{chapters@#1}%
    {\PackageError{chapters}{Sheet #1 is not listed in the index}%
      {List it in 20261125T012000--grh-workbook__guide.org\MessageBreak
       and run scripts/make-chapters.py.}}%
    {\edef\thechap{\@nameuse{chapters@#1}}}}
\makeatother
\chapnumber{20261126T050500}
\renewcommand{\thesection}{\thechap.\arabic{section}}

\theoremstyle{plain}
\newtheorem{theorem}{Theorem}[section]
\newtheorem{corollary}[theorem]{Corollary}
\theoremstyle{definition}
\newtheorem{problem}{Problem}[section]
\newtheorem{definition}[theorem]{Definition}

\newcommand{\Q}{\mathbb{Q}}
\newcommand{\C}{\mathbb{C}}
\newcommand{\Oc}{\mathcal{O}}
\newcommand{\Mcal}{\mathcal{M}}
\newcommand{\Rc}{\mathcal{R}}
\newcommand{\TM}{\mathbb{T}^{\mathcal{M}}}
\newcommand{\tr}{\operatorname{tr}}
\newcommand{\tl}[1]{\texttt{#1}}

\setlength{\emergencystretch}{3em}

\title{Equal multiplicities}
\author{}
\date{November 26, 2026}

\begin{document}
\maketitle
\refstepcounter{chap}\label{chap}

\paragraph{Prerequisites.}
Linear algebra over \(\C\): traces, characteristic polynomials and Newton's identities \cite[Ch.~I, \S2]{Mac}.  Modular forms of weight two on \(\Gamma_0(N)\) and their Hecke operators \cite[Ch.~5]{DS}.  From Chapter~\ref{ho-chap}, \emph{Hecke operators and theta lifts}: Section~\ref{ho-sec:brandt}, with Problems~\ref{ho-pr:brandt} and~\ref{ho-pr:basis}.  From Chapter~\ref{ob-chap}, \emph{Spaces, forms and twists}: the chosen places \(\Mcal\), the group \(\bar G\), good families of labels and the idempotent \(e_0\) (Section~\ref{ob-sec:chosen}); the gap \(\Delta\) and the target parameter \(\varsigma_\tau\) (Section~\ref{ob-sec:weights}, Problem~\ref{ob-pr:parameter}); the spaces \(Y_0\), \(W_j\) and their signs (Sections~\ref{ob-sec:spaces} and~\ref{ob-sec:forms}, Problem~\ref{ob-pr:signs-W}); principal and good forms \(H(V) \supseteq H_g(V)\) and the target \(T = H_g(Y_0)\) of dimension \(d_T\); tame characters (Section~\ref{ob-sec:twists}).  From Chapter~\ref{da-chap}: \(r = 3^{e'}\) is odd, \(x^5 \le r\), \(m\) is even, and the weights \(\kappa(\tau)\) are integers.  From Chapter~\ref{c1-chap}: the Rallis operators \(\Rc_j(\xi; V)\) (Definition~\ref{c1-def:rallis-operator}).  Inner products are conjugate-linear in the first variable.  Problems marked ``by machine'' use Sage \cite{Sage}.

\medskip
The single lifts of Chapter~\ref{dch-chap}, \emph{The determinant of a chain}, start from the spaces \(W_j\), which have the rank of \(Y_0\) but other signs at the real places.  How can a determinant on \(H_g(W_j)\) be compared with one on \(T = H_g(Y_0)\), when no map between the two spaces is given?  In the proof one Hecke algebra acts on both spaces with equal traces.  The equality gives their dimensions and characteristic polynomials and, with the decompositions into joint eigenspaces, the multiplicities of the Hecke eigensystems.  This sheet derives the consequences, and states the equality (Theorem~\ref{thm:equal-traces}) with the comparison of trace formulas behind it.

\section{Hecke operators of a finite group}\label{sec:finite}

Here \(K\) is a subgroup (the field \(\Q(\sqrt{-d})\) does not occur).  Let the symmetric group \(S_3\) act on \(\C^3\) by \((\sigma v)_i = v_{\sigma^{-1}(i)}\), let \(K = \{1, (12)\}\) and \(g = (23)\).

\begin{problem}\label{pr:s3-coset}
The double coset \(KgK\) has \(4\) elements.  For a \(K\)-fixed vector \(v = (v_1, v_1, v_3)\), the average of \(\sigma v\) over \(\sigma \in KgK\) is \(\bigl(\tfrac{v_1+v_3}2, \tfrac{v_1+v_3}2, v_1\bigr)\); in the coordinates \((v_1, v_3)\) of the \(K\)-fixed vectors it is \((v_1, v_3) \mapsto \bigl(\tfrac{v_1+v_3}2, v_1\bigr)\).
\end{problem}

\begin{problem}\label{pr:compress}
Let a group \(G\) act on a complex inner product space \(U\) by unitary operators, let \(K \subseteq G\) be a subgroup and \(g \in G\) be such that \(K_g := K \cap gKg^{-1}\) has finite index in \(K\).  On the \(K\)-fixed vectors \(U^K\) the operator \(\mathrm{av}_K(g) := [K : K_g]^{-1} \sum_{\delta K_g \in K/K_g} \delta g\) does not depend on the representatives, maps \(U^K\) to itself, and satisfies \(\langle z, gy \rangle = \langle z, \mathrm{av}_K(g) y\rangle\) for all \(z, y \in U^K\).  If \(K\) is finite, \(\mathrm{av}_K(g) y\) is the average of \(\sigma y\) over \(\sigma \in KgK\).
\end{problem}

\begin{problem}\label{pr:s3-traces}
The operator of Problem~\ref{pr:s3-coset} has traces \(2, \tfrac12, \tfrac54, \tfrac78\) in its powers \(0, 1, 2, 3\), characteristic polynomial \((X - 1)(X + \tfrac12)\) and determinant \(-\tfrac12\).
\end{problem}

\begin{problem}\label{pr:trace-charpoly}
Let a monoid act linearly, with its unit acting as the identity and products as composites, on finite-dimensional complex spaces \(U\) and \(U'\), and suppose \(\tr_U h = \tr_{U'} h\) for every element \(h\).  Then \(\dim U = \dim U'\), and every element \(h\) has the same characteristic polynomial, hence the same determinant, on \(U\) and on \(U'\).
\end{problem}

\begin{problem}\label{pr:trace-check}
The matrices \(\operatorname{diag}(1, 2)\) and \(\bigl(\begin{smallmatrix} 0 & -2 \\ 1 & 3 \end{smallmatrix}\bigr)\) both have traces \(3, 5, 9\) in their powers \(1, 2, 3\), and \(\det(h + 1) = 6\) for both.
\end{problem}

\begin{problem}\label{pr:trace-mult}
In Problem~\ref{pr:trace-charpoly}, suppose moreover that the monoid is commutative and acts diagonalizably on \(U\) and on \(U'\).  Then for every function \(\lambda\) from the monoid to \(\C\) the joint eigenspaces \(\{y : hy = \lambda(h) y \text{ for all } h\}\) in \(U\) and in \(U'\) have the same dimension.
\end{problem}

\section{Brandt matrices: a model}\label{sec:brandt}

Where do two spaces with the same Hecke traces come from?  The classical case compares a definite quaternion algebra with \(\mathrm{GL}_2\).  Let \(B\), \(\Oc_1\), \(\Oc_2\), the Brandt matrices \(B(\ell)\) for primes \(\ell \ne 11\), and the cusp form \(f\) be as in Section~\ref{ho-sec:brandt} and Problem~\ref{ho-pr:basis}; \(B(\ell)\) acts on functions \(y\) on the two classes by \((B(\ell) y)(i) = \sum_j B_{ij}(\ell)\, y(j)\).  The space of cusp forms of weight \(2\) on \(\Gamma_0(11)\) is \(\C f\) \cite[\S3.9]{DS}, and \(T_\ell f = a_\ell(f) f\).

\begin{problem}\label{pr:level11}
(By machine.)  For \(\ell = 2, 3, 5, 7, 13\): \(\tr B(\ell) = 1, 3, 7, 6, 18\), and \(\tr B(\ell) - (\ell + 1) = a_\ell(f) = -2, -1, 1, -2, 4\).
\end{problem}

\begin{problem}\label{pr:remove-constants}
(a) If an idempotent \(e\) on a finite-dimensional space \(U\) commutes with an operator \(h\), then \(eU\) is stable under \(h\) and \(\tr(h \mid eU) = \tr_U(he)\).  (b) For a prime \(\ell \ne 11\), the constant function is an eigenvector of \(B(\ell)\) with eigenvalue \(\ell + 1\), and the line \(U_0 := \{y : y(1)/\lvert\Oc_1^\times\rvert + y(2)/\lvert\Oc_2^\times\rvert = 0\}\) is stable under \(B(\ell)\); hence \(\tr(B(\ell) \mid U_0) = \tr B(\ell) - (\ell + 1)\).
\end{problem}

For the five primes of Problem~\ref{pr:level11}, Problems~\ref{pr:level11} and~\ref{pr:remove-constants} give \(\tr(B(\ell) \mid U_0) = a_\ell(f)\): once the constants are removed, the Brandt matrices have the traces of the Hecke operators on cusp forms.  The equality for every prime \(\ell \ne 11\) is the level-\(11\) case of the correspondence of Eichler and of Jacquet and Langlands \cite{Ei73}, \cite[\S16]{JL}; \emph{Status}: classical.  It is a model only: the group of units of \(B\) and \(\mathrm{GL}_2\) differ at \(11\) as well as at infinity, whereas the proof compares unitary groups that differ only at the real places.

\begin{problem}\label{pr:level37}
(By machine.)  The quaternion algebra over \(\Q\) ramified at \(37\) and \(\infty\) has three classes of left ideals of a maximal order; its Brandt matrix \(B(2)\) has characteristic polynomial \((X - 3)X(X + 2)\), and \(T_2\) on the cusp forms of weight \(2\) on \(\Gamma_0(37)\) has characteristic polynomial \(X(X + 2)\).
\end{problem}

At level \(37\) the factor of the constants is \(X - 3\), and what remains is the characteristic polynomial of \(T_2\) on cusp forms.

\section{Equal traces}\label{sec:traces}

Let \(V\) be a framed space of rank \(r\) over \(E\) with numbered signs \(s = (s(0), \dots, s(r-1))\) (in this sheet \(s\) is a system of signs and, below, \(a\) an element of a group algebra; the numbers \(s\) and \(a\) of the data do not occur), such as \(Y_0\) or \(W_j\) (Section~\ref{ob-sec:forms}).  Through the frames, all such \(V\) have the same finite adelic unitary group \(G_f\), that of the identity form; it acts on forms on \(V\) by right translation.  Give \(H(V)\) and its subspaces the \emph{Petersson inner product} \(\langle y', y\rangle := \int \overline{y'(g)}\, y(g)\, dg\), the integral over the quotient of the adelic unitary group of \(V\) by its rational points, for the quotient measure fixed in the proof.

\begin{definition}\label{def:gap}
The numbered signs \(s\) of \(V\) satisfy the \emph{gap hypothesis} if at every infinite place \(\tau \in \Phi\) the entries \(2\varsigma_\tau(0), \dots, 2\varsigma_\tau(r-1)\) are distinct, are congruent to \(r - 1\) modulo \(2\), and satisfy \(\lvert 2\varsigma_\tau(i) - 2\varsigma_\tau(i')\rvert \ge 2\Delta\) whenever \(s(i) \ne s(i')\).
\end{definition}

Under the gap hypothesis the proof shows that \(H(V)\) is finite-dimensional; traces below are taken under this hypothesis.

\begin{problem}\label{pr:gap-Y-W}
The gap hypothesis holds for \(Y_0\) and for \(W_j\), \(1 \le j \le m\).
\end{problem}

\begin{definition}\label{def:hecke-away}
In this sheet \(G_f^{\Mcal}\) is the subgroup of the \(g \in G_f\) whose component at every \(w \in \Mcal\) is \(1\); the component at \(w\) determines that at \(\bar w\), so these components are \(1\) as well.  An element \(a\) of the group algebra \(\C[G_f^{\Mcal}]\) is \emph{spherical} if, in every representation of \(G_f\), it maps the vectors fixed by the principal level to vectors fixed by the principal level.  A spherical \(a\) acts on \(H(V)\) for every framed \(V\) of rank \(r\); the \emph{Hecke algebra away from \(\Mcal\)}, \(\TM\), is the algebra of the families of operators on the spaces \(H(V)\), one for each \(V\), given in this way by one spherical element.  For a finite place \(w\) and \(i \in \mathbb N\), let \(h_{w,i} \in \TM\) be the Hecke operator (the sum of the translates by representatives of its single cosets) of the \(\mathrm{GL}_r(\Oc_w)\)-double coset of \(\operatorname{diag}(\varpi_w 1_i, 1_{r-i})\) in \(\mathrm{GL}_r(E_w)\), with \(\Oc_w\) the local ring of integers and \(\varpi_w\) a uniformizer, if \(w\) is split, \(w, \bar w \notin \Mcal\) and \(i \le r\); in all other cases \(h_{w,i} := 0\).
\end{definition}

In this sheet an \emph{eigensystem} is a family \(\lambda = (\lambda_{w,i})\) of complex numbers indexed by the finite places \(w\) and \(i \in \mathbb N\), and \(H(V)_\lambda := \{y \in H(V) : h_{w,i} y = \lambda_{w,i} y \text{ for all } w, i\}\).  The proof lets \(\bar G\) act on \(H(V)\) through integral lifts whose components outside \(\Mcal \cup \overline{\Mcal}\) are \(1\).

\begin{problem}\label{pr:commute}
An element of \(G_f^{\Mcal}\) commutes with every element of \(G_f\) whose components outside \(\Mcal \cup \overline{\Mcal}\) are \(1\).  Hence every \(h \in \TM\) commutes with the action of \(\bar G\), and so with \(e_0\), on \(H(V)\), and maps \(H_g(V)\) to itself.
\end{problem}

For \(V\) satisfying the gap hypothesis and \(h \in \TM\), write \(\tr(h \mid H_g(V))\) for the trace of the restriction of \(h\) to \(H_g(V)\).  The proof establishes the following.

\begin{theorem}\label{thm:equal-traces}
Let \(V\) and \(V'\) be framed spaces of rank \(r\) over \(E\) whose numbered signs satisfy the gap hypothesis.  Then every \(h \in \TM\) has \(\tr(h \mid H_g(V)) = \tr(h \mid H_g(V'))\).
\end{theorem}

Theorem~\ref{thm:equal-traces} is not proved here.  The proof deduces it from a comparison of simple trace formulas.  Let \(E\) be a CM field with a fixed CM type, containing the \(5040\)th roots of unity and having an even number of infinite places.  Let \(V\), \(V'\) be framed spaces of rank \(n \ge 2\) over \(E\) with numbered signs \(s\), \(s'\), and \(G_f\) their common finite adelic group.  For each infinite place \(\tau\) let \(l_\tau \in \mathbb Z^n\) have distinct entries congruent to \(n - 1\) modulo \(2\), with \(\lvert l_\tau(i) - l_\tau(i')\rvert \ge 2(500 n^3 + 10 n^5)\) whenever \(s(i) \ne s(i')\) or \(s'(i) \ne s'(i')\).  Let \(K_f\) be a subgroup of the group of integral elements of \(G_f\) that contains a neighbourhood of the identity, and let \(a \in \C[G_f]\) map \(K_f\)-fixed vectors to \(K_f\)-fixed vectors in every representation of \(G_f\).  Suppose that the test function \(a \star 1_{K_f}\) is \emph{cuspidal} at a split place \(w\), with residue field \(k_w\): for every nonempty proper set \(I\) of coordinates, with \(\mathrm{pr}_I\) the diagonal projection onto it, every \(\bar g \in \mathrm{GL}_n(k_w)\), every family \((u_A)_{A \in M_n(k_w)}\) in \(\mathrm{GL}_n(\Oc_w)\) with \(u_A\) reducing to \((1 + \mathrm{pr}_I A (1 - \mathrm{pr}_I))\bar g\), and every \(g \in G_f\), the sum of \((a \star 1_{K_f})(g\, u_A)\) over \(A \in M_n(k_w)\) is zero, \(u_A\) being placed in \(G_f\) at \(w\).  Then operators acting as \(a\) on the \(K_f\)-fixed forms of parameter \(l\) on \(V\) and on \(V'\) (defined like the principal forms, with \(l_\tau\) for \(2\varsigma_\tau\) and \(K_f\) for the principal level) have the same trace.  For the classical arguments, see \cite[Thms.~5.4--5.5]{Lab}, \cite[Ch.~1, Lemmas~2.4--2.6]{AC} and, for a cuspidal component, \cite[Prop.~1.10.5]{HL}; \emph{Status}: the method is classical; the proof establishes the comparison in this form, with integrable idempotent kernels of discrete series in place of pseudo-coefficients.

The data meet these conditions: the proof shows that \(E\) contains the \(5040\)th roots of unity and has an even number of infinite places, and \(\Delta = 2r^8 \ge 500 r^3 + 10 r^5\) and \(r \ge 2\) because \(r \ge x^5 \ge 9\), \(x\) being large.  The idempotent \(e_0\) is the sum of the orthogonal idempotents \(e_c := \bigotimes_w e_{c_w}\) over the good families \(c\), so by Problem~\ref{pr:remove-constants}(a) the trace of \(h \in \TM\) on \(H_g(V)\) is \(\sum_c \tr_{H(V)}(h e_c)\).  The proof shows that \(h e_c\) acts on \(H(V)\) for every \(V\) through one element of \(\C[G_f]\) that keeps the principal-level vectors, and that its test function is cuspidal at a chosen place \(w\) where \(c_w\) is primitive (irreducible with a root of order \((\#k_w)^r - 1\)); every good family has such a place (Section~\ref{ob-sec:chosen}), and chosen places are split.  The comparison applies to each term, with \(K_f\) the principal level and \(l_\tau = 2\varsigma_\tau\).

\section{Dimensions, eigenspaces and determinants}\label{sec:corollary}

\begin{problem}\label{pr:dim}
\(\dim H_g(W_j) = d_T\) for \(1 \le j \le m\).
\end{problem}

When \(H(V)\) is finite-dimensional, the proof shows that it is the direct sum of the spaces \(H(V)_\lambda\) (the operators \(h_{w,i}\) are normal and commute).

\begin{problem}\label{pr:eigen-proj}
Let \(\lambda\) be an eigensystem and fix finitely many eigensystems \(\lambda' \ne \lambda\).  A polynomial in finitely many of the \(h_{w,i}\), an element of \(\TM\), acts on \(H(V)\), for every framed \(V\) of rank \(r\), as the identity on \(H(V)_\lambda\) and as zero on each of these \(H(V)_{\lambda'}\).
\end{problem}

\begin{problem}\label{pr:eigen}
For every eigensystem \(\lambda\) and \(1 \le j \le m\), the spaces \(T \cap H(Y_0)_\lambda\) and \(H_g(W_j) \cap H(W_j)_\lambda\) have the same dimension.
\end{problem}

Let \(\xi\) be a tame character and \(j \in \mathbb N\).  By Definition~\ref{c1-def:rallis-operator} and the formula for its matrix coefficients proved with it, for every framed \(V\) of rank \(r\) and all \(z, y \in H(V)\),
\[
  \langle z, \Rc_j(\xi; V)\, y \rangle = \sum_\gamma \mathrm{mc}(\gamma)\, \langle z, \gamma y\rangle ,
\]
a sum, convergent in any order, over the countable family of representatives \(\gamma \in G_f\) of that definition, the same family for every \(V\), whose components at the places of \(\Mcal\) are \(1\), with the coefficients \(\mathrm{mc}(\gamma)\) of Chapter~\ref{c1-chap}, which depend on \(\xi\), \(j\) and \(\gamma\) but not on \(V\); here \(\gamma y := y(\cdot\,\gamma)\).  For each \(\gamma \in G_f^{\Mcal}\) the proof supplies an element \(\mathrm{av}(\gamma) \in \TM\), an average of translates as in Problem~\ref{pr:compress}, such that \(\langle z, \mathrm{av}(\gamma) y\rangle = \langle z, \gamma y\rangle\) for every framed \(V\) of rank \(r\) and all \(z, y \in H(V)\).

\begin{problem}\label{pr:closed}
Let \(U\) and \(U'\) be finite-dimensional complex inner product spaces, and let pairs \((h_i, h'_i) \in \operatorname{End} U \times \operatorname{End} U'\), for \(i\) in a countable set, lie in one complex linear subspace.  Let \(h_\infty \in \operatorname{End} U\) and \(h'_\infty \in \operatorname{End} U'\) be such that \(\sum_i \langle z, h_i y\rangle = \langle z, h_\infty y\rangle\) for all \(z, y \in U\) and \(\sum_i \langle z', h'_i y'\rangle = \langle z', h'_\infty y'\rangle\) for all \(z', y' \in U'\), the sums converging in any order.  Then \((h_\infty, h'_\infty)\) lies in that subspace.  In particular, let a complex algebra act on \(U\) and on \(U'\) through complex-linear algebra homomorphisms to \(\operatorname{End} U\) and \(\operatorname{End} U'\); if each \((h_i, h'_i)\) is the action of one element, then one element acts as \(h_\infty\) on \(U\) and as \(h'_\infty\) on \(U'\).
\end{problem}

\begin{problem}\label{pr:rallis-hecke}
Let \(V\) and \(V'\) be framed spaces of rank \(r\) satisfying the gap hypothesis, \(\xi\) a tame character and \(j \in \mathbb N\).  One element of \(\TM\) acts as \(\Rc_j(\xi; V)\) on \(H(V)\) and as \(\Rc_j(\xi; V')\) on \(H(V')\).  In particular \(\Rc_j(\xi; V)\) maps \(H_g(V)\) to itself.
\end{problem}

\begin{problem}\label{pr:det}
For \(1 \le j \le m\) and every tame character \(\xi\), \(\det\bigl(\Rc_j(\xi; W_j) \mid H_g(W_j)\bigr) = \det\bigl(\Rc_j(\xi; Y_0) \mid T\bigr)\).
\end{problem}

Problems~\ref{pr:dim}, \ref{pr:eigen}, \ref{pr:rallis-hecke} and~\ref{pr:det} prove the following.

\begin{corollary}\label{cor:equal}
Let \(1 \le j \le m\).  Then \(\dim H_g(W_j) = d_T\); for every eigensystem \(\lambda\) the \(\lambda\)-eigenspaces of the \(h_{w,i}\) in \(T\) and in \(H_g(W_j)\) have the same dimension; and for every tame character \(\xi\), \(\Rc_j(\xi; Y_0)\) maps \(T\) to itself, \(\Rc_j(\xi; W_j)\) maps \(H_g(W_j)\) to itself, and \(\det\bigl(\Rc_j(\xi; W_j) \mid H_g(W_j)\bigr) = \det\bigl(\Rc_j(\xi; Y_0) \mid T\bigr)\).
\end{corollary}

The problems deduce Corollary~\ref{cor:equal} from Theorem~\ref{thm:equal-traces}, which rests on the comparison of trace formulas taken as given, and from facts read from the proof: finite-dimensionality under the gap hypothesis, the decomposition into joint eigenspaces, the series for the matrix coefficients of \(\Rc_j(\xi; V)\), and the operators \(\mathrm{av}(\gamma)\).  Chapter~\ref{dch-chap} combines the corollary with Theorem~\ref{c1-thm:one-lift}, which writes \(\hat u_{\xi,j}^*\hat u_{\xi,j}\) as a multiple of \(\Rc_j(\xi; W_j)\), to express the volume of a single lift by a determinant on \(T\) (Problem~\ref{dch-pr:one-lift}).

\paragraph{In the Lean formalization.}
\begin{sloppypar}
In the Lean proof \cite{GL} (\tl{grh/src/}, namespace \tl{New.Data} unless written in full; EM and HE are the files \tl{EqualMultiplicities.lean} and \tl{HeckeEigensystems.lean} in \tl{New/Spanning/Multiplicities/}), Theorem~\ref{thm:equal-traces} is \tl{trace\_heckeOnGood\_eq} (EM).  Characteristic polynomials: \tl{charpoly\_heckeOnGood\_eq}.  The dimension assertion of Corollary~\ref{cor:equal} is \tl{targetDim\_eq\_finrank\_goodForms\_spaceW}, for \(Y_0, W_j\) with \(1 \le j \le m\); the multiplicity assertion specializes \tl{finrank\_goodForms\_inf\_eigenPart\_eq}, stated for two spaces with the gap hypothesis; \tl{rallisOnGood} restricts the Rallis operator to good forms on one, for any tame character and \(j \in \mathbb N\), and \tl{det\_rallisOnGood\_eq} equates the determinants on two, for the same character and \(j\).  All of these are in EM.  Definition~\ref{def:gap} is \tl{New.HCParam.HasGap} for \tl{paramW} and \tl{gap}; the gap for \(Y_0\), \(W_j\) is \tl{hasGap\_signsY\_paramW}, \tl{hasGap\_paramW}, and finite-dimensionality the instance after \tl{isCuspKernel\_typeKernel\_paramW}, in \tl{New/UpperBound/Objects/Rows.lean}.  Definition~\ref{def:hecke-away} uses \tl{offLevelGroup}, \tl{sphericalAlgebra}, \tl{heckeAlgebra} and \tl{heckeOperator} (HE).  In HE the restriction to good forms is \tl{heckeOnGood}, the eigenspaces are \tl{eigenPart}, and their direct sum is given by \tl{iSup\_eigenPart\_eq\_top} (spanning) and \tl{iSupIndep\_eigenPart} (directness).  The action of \(\bar G\) is given by \tl{coe\_principalRep\_apply} (HE) with \tl{New.liftResidue\_mem\_supportedOn}, and commutation with it follows from \tl{commute\_heckeAlgebra\_principalRep} (HE).  The proof of \tl{exists\_hecke\_proj\_eigenPart} (HE) constructs the polynomial of Problem~\ref{pr:eigen-proj}.  At split places away from the level, \tl{New.heckeSplit\_adjoint} identifies the adjoint on principal forms with the operator at the conjugate place; these operators commute by \tl{New.heckeSplit\_comm}, which gives normality as well.  The comparison of trace formulas is \tl{New.trace\_eq\_of\_isCuspidalAt} (\tl{Lib/Adelic/TraceFormula/CuspidalComparison.lean}), with \tl{New.IsLevel} and \tl{New.IsCuspidalAt}; the reduction to it is \tl{trace\_heckeOnGood\_eq\_sum}, \tl{exists\_isPrimitivePolynomial} and \tl{trace\_mul\_tupleIdem\_eq} (EM), and the field conditions on the data follow from the instance of \tl{New.HasPrimitiveRoot5040} for \(E\) in \tl{New/Derivatives/Field/PrimeDecomposition.lean} and from \tl{f\_even}.  The series for \(\Rc_j(\xi; V)\) is \tl{hasSum\_inner\_rallisOperator} (\tl{New/Determinant/Norms/AdjointThetaLift.lean}), its representatives and their countability \tl{New.SiegelWeil.cosetRep} and \tl{New.SiegelWeil.countable\_quotient\_localLevel}; the operators \(\mathrm{av}(\gamma)\) and the common operator of Problem~\ref{pr:rallis-hecke} are \tl{exists\_heckeAlgebra\_inner\_eq\_inner\_translate} and \tl{exists\_hecke\_eq\_rallisOperator} (EM).
\end{sloppypar}

\section*{Hints}

\subsection*{\S\ref{sec:finite}}
\begin{description}
\item[\ref{pr:s3-coset}] \(KgK = Kg \cup K g (12)\); list the four permutations.
\item[\ref{pr:compress}] For \(h \in K_g\), \(g^{-1}hg \in K\) fixes \(y\).  Move \(\delta\) to the other side of the inner product.
\item[\ref{pr:s3-traces}] Write the operator as a \(2 \times 2\) matrix; the traces of powers are sums of powers of its eigenvalues.
\item[\ref{pr:trace-charpoly}] Take \(h = 1\), then the traces of the powers of \(h\), and Newton's identities.
\item[\ref{pr:trace-mult}] Only finitely many \(\lambda\) occur on \(U \oplus U'\).
\end{description}

\subsection*{\S\ref{sec:brandt}}
\begin{description}
\item[\ref{pr:level11}] Sage: the functions \tl{BrandtModule} (with its Hecke matrices) and \tl{EllipticCurve} (for the curve \(X_0(11)\), whose traces of Frobenius are the \(a_\ell(f)\)).  For \(\ell = 2, 3\), Problems~\ref{ho-pr:brandt} and~\ref{ho-pr:basis}.
\item[\ref{pr:remove-constants}] (b) The row sums of \(B(\ell)\) and the relation between \(B_{ij}(\ell)\) and \(B_{ji}(\ell)\) quoted in Section~\ref{ho-sec:brandt}; the constants are not in \(U_0\).
\item[\ref{pr:level37}] Sage: the functions \tl{BrandtModule} and \tl{CuspForms}.
\end{description}

\subsection*{\S\ref{sec:traces}}
\begin{description}
\item[\ref{pr:gap-Y-W}] Problem~\ref{ob-pr:parameter} and its hint.  For \(Y_0\) the signs alternate along the entries after the semicolon.  For \(W_j\) they change from \(+1\) to \(-1\) at index \(k + j - 1\) if \(j < m\), and all are \(+1\) if \(j = m\).
\item[\ref{pr:commute}] Compare components place by place; \(e_0 \in \C[\bar G]\).
\end{description}

\subsection*{\S\ref{sec:corollary}}
\begin{description}
\item[\ref{pr:eigen-proj}] For each \(\lambda'\) choose \((w, i)\) with \(\lambda_{w,i} \ne \lambda'_{w,i}\); no commutation is needed.
\item[\ref{pr:eigen}] Only finitely many eigensystems occur in \(H(Y_0)\) or \(H(W_j)\).  Compare with Problem~\ref{pr:trace-mult}.
\item[\ref{pr:closed}] Finite-dimensional subspaces are closed.
\item[\ref{pr:rallis-hecke}] Take \(K\) to be the group of integral elements of \(G_f^{\Mcal}\).  It fixes the principal forms, and \(K_\gamma := K \cap \gamma K \gamma^{-1}\) has finite index in \(K\).  For a choice of representatives of \(K/K_\gamma\), the proof shows that the normalized sum of the \(\delta\gamma\) in \(\C[G_f^{\Mcal}]\) is spherical; it induces the family \(\mathrm{av}(\gamma)\), whose action on principal forms is the average of Problem~\ref{pr:compress}.  For the series, Problem~\ref{pr:closed} with the elements \(\mathrm{mc}(\gamma)\,\mathrm{av}(\gamma)\).
\end{description}

{\small
\begin{thebibliography}{Sage}
\bibitem[AC]{AC} J.~Arthur and L.~Clozel, \emph{Simple Algebras, Base Change, and the Advanced Theory of the Trace Formula}, Ann.\ of Math.\ Stud.\ 120, Princeton Univ.\ Press, 1989.
\bibitem[DS]{DS} F.~Diamond and J.~Shurman, \emph{A First Course in Modular Forms}, Grad.\ Texts in Math.\ 228, Springer, 2005.
\bibitem[Ei73]{Ei73} M.~Eichler, \emph{The basis problem for modular forms and the traces of the Hecke operators}, in \emph{Modular Functions of One Variable I}, Lecture Notes in Math.\ 320, Springer, 1973, 75--151.
\bibitem[GL]{GL} L.~Sernau, \emph{GRH for Dirichlet \(L\)-functions, in Lean~4}, a shorter rewriting of the Lean proof, text by Claude, \url{https://github.com/lsernau/grh-lean}, commit \texttt{72fead1}, readable edition \texttt{grh/}.
\bibitem[HL]{HL} M.~Harris and J.-P.~Labesse, \emph{Conditional base change for unitary groups}, Asian J.\ Math.\ \textbf{8} (2004), 653--683.
\bibitem[JL]{JL} H.~Jacquet and R.~P.~Langlands, \emph{Automorphic Forms on \(\mathrm{GL}(2)\)}, Lecture Notes in Math.\ 114, Springer, 1970.
\bibitem[Lab]{Lab} J.-P.~Labesse, \emph{Changement de base CM et s\'eries discr\`etes}, in \emph{On the Stabilization of the Trace Formula}, Stab.\ Trace Formula Shimura Var.\ Arith.\ Appl.\ 1, Int.\ Press, 2011, 429--470.
\bibitem[Mac]{Mac} I.~G.~Macdonald, \emph{Symmetric Functions and Hall Polynomials}, 2nd ed., Oxford University Press, 1995.
\bibitem[Sage]{Sage} The Sage Developers, \emph{SageMath, the Sage Mathematics Software System}, \url{https://www.sagemath.org}.
\end{thebibliography}
}

\end{document}
